\documentclass{article}
\usepackage{amsmath}
\usepackage{xcolor}
\begin{document}

E.D. DE PRIMER ORDEN 

En este tipo de E.D. solo se cuenta con la primera derivada presente en la E.D.

\textbf{Representacion}

Se puede tener la E.D. en su forma \colorbox[RGB]{248,231,28}{\textbf{Estandar:}}


\begin{gather*}
\frac{dy}{dx} = f(x,y)  \rightarrow  A \\
\frac{dy}{dx} = xe^{y}
\end{gather*}


Se puede tener asi mismo su representacion de forma \colorbox[RGB]{248,231,28}{\textbf{diferencial}}


\begin{gather*}
M(x,y)dx + N(x,y)dy = 0  \rightarrow  B \\
(x^{3}+5y)dx+(e^{2x}-ln(y))dy = 0
\end{gather*}


Es posible readecuar la representacion algebraida de una E.D. para converirtirla

a cualquiera de sus formas

Por ejmplo:  


\begin{gather*}
5x+3y-\frac{dy}{dx} = 0 \\
\text{para llevar a la forma A} \\
\frac{dy}{dx} = 5x+3y \\
\text{es posible llevar la E.D. a su forma B} \\
\text{podriamos multiplicar la expresion por dx} \\
\frac{dy}{\text{\colorbox[RGB]{248,231,28}{$dx$}}}*\text{\colorbox[RGB]{248,231,28}{$dx$}} = (5x+3y)dx \\
dy = (5x+3y)dx \\
M(x,y)dx + N(x,y)dy = 0 \\
(5x+3y)dx-dy = 0   \rightarrow  \\
M(x,y) =(5x+3y); N(x,y) = -1
\end{gather*}




Clasificacion

\textbf{E.D. de Variable Separable}

En este tipo de E.D. es posible aplicar directamente la integracion para resolver la misma, ya que se tiene la certeza de que cada variable esta separada y sujeta a su respectiva diferencial




\begin{gather*}
\frac{dy}{dx} = f(x,y)  \rightarrow \text{ donde es posible tener una empresion } \\
\text{donde se pueda despejar individualmente cada variable} \\
\frac{dy}{dx} = A(x)*B(y) o A(x)/B(y) o B(x)/A(y)\text{   o las posibles combinaciones de grupos} \\
\text{* supongramos que tenemos el primer caso} \\
\frac{dy}{dx} = A(x)*B(y)   \rightarrow dy =\text{\colorbox[RGB]{248,231,28}{$A$}}\text{\colorbox[RGB]{248,231,28}{$($}}\text{\colorbox[RGB]{248,231,28}{$x$}}\text{\colorbox[RGB]{248,231,28}{$)$}}*B(y)\text{\colorbox[RGB]{248,231,28}{$dx$}}  \rightarrow \frac{1}{B(y)}dy =\text{\colorbox[RGB]{248,231,28}{$A$}}\text{\colorbox[RGB]{248,231,28}{$($}}\text{\colorbox[RGB]{248,231,28}{$x$}}\text{\colorbox[RGB]{248,231,28}{$)$}}\text{\colorbox[RGB]{248,231,28}{$dx$}}  \\
\int \frac{1}{B(y)}dy =\text{\colorbox[RGB]{248,231,28}{$\int $}}A(x)dx  \\
o \\
\text{partiendo de la forma B} \\
M(x,y)dx + N(x,y)dy = 0   \rightarrow \text{ donde} \\
M(x,y)  \rightarrow  A(x) \\
N(x,y) \rightarrow B(y) \\
A(x)dx + B(y)dy = 0 \\
\text{se puede aplicar integral directa} \\
\int A(x)dx + \int B(y)dy =c
\end{gather*}




\textbf{E.D Homogeneas} $\rightarrow$  no se pueden separar directamente, pero si se puede hacer un cambio de variable para resolver la misma convirtiendola a una E.D. de Variable Separable




\begin{gather*}
\frac{dy}{dx} = f(x,y) \\
\text{Si se incorporar el operador lambda }\text{\colorbox[RGB]{248,231,28}{$𝜆$}}\text{ a cada variable } \\
\text{y se puede nuevamente reproducir la  expresion original,} \\
\text{entonces la E.D. es Homogenea} \\
\vspace{0.5em} \\
\frac{dy}{dx} = f(x,y) = \frac{dy}{dx} = f(\text{\colorbox[RGB]{248,231,28}{$𝜆$}}x,\text{\colorbox[RGB]{248,231,28}{$𝜆$}}y) =f(x,y)   \\
\text{entonces se hace un cambio de variable} \\
y = xv    o      x =yu \\
\frac{dy}{dx} = v+x\frac{dv}{dx}   o    \frac{dx}{dy} = u+y\frac{du}{dy} \\
\text{una vez acomodada algebraicamente, } \\
\text{se puede resolver por Variable Separable}
\end{gather*}


Lineales


\begin{gather*}
\frac{dy}{dx}+P(x)y =Q(x)  
\end{gather*}


E.D. de Bernoulli $\rightarrow$  es una variacion a una E.D. Lineal


\begin{gather*}
\frac{dy}{dx}+P(x)y =Q(x)\text{\colorbox[RGB]{248,231,28}{$y$}}\text{\colorbox[RGB]{248,231,28}{$^{n}$}}
\end{gather*}




E.D Exactas


\begin{gather*}
M(x,y)dx + N(x,y)dy = 0 \\
\text{para verificar que una E.D. es exacta,} \\
\text{se utiliza prueba de exactitud.} \\
\frac{\partial M(x,y)}{\partial y} = \frac{\partial N(x,y)}{\partial x}
\end{gather*}









\end{document}
